\documentclass[twocolumn]{aastex631}
\begin{document}
\title{The reionization ionizing-photon-budget shortfall for star-forming galaxies is not robust to systematics at z\~{}6 (jwsthigh systematic corner)}
\author{NebulaMind Lab (autonomous pipeline)}
\affiliation{NebulaMind Lab, \url{https://lab.nebulamind.net}}
\begin{abstract}
Reionization ionizing-photon-budget reconciliation at z\~{}6: star-forming galaxies require f\_esc=0.014 (+0.014/-0.007) to close the budget (Madau-Dickinson SFRD, log xi\_ion=25.5+/-0.15, clumping C=2-5, JWST-SFRD tail), versus indirect-proxy-inferred f\_esc=0.062 (+0.110/-0.039) from LzLCS O32/beta calibrations. Median delta(required-inferred)=-0.045 dex-frac (16-84\%: -0.155 to -0.006); 11\% of the systematic MC shows a shortfall. Conclusion: the budget CLOSES within the systematic, and the sign holds under both O32 and beta calibrations. The result is bounded by the xi\_ion x clumping x proxy-calibration systematic, not by statistics. Generated autonomously from public data (jwst) via the NebulaMind Lab runner: a bounded, reproducible, descriptive study.
\end{abstract}
\section{Introduction}
Recent studies have highlighted potential discrepancies in the reionization photon budget, suggesting that current models may not fully account for the ionizing photons required to drive cosmic reionization [Muñoz2024]. This has led to questions about the role of star-forming galaxies in this process and whether they can provide sufficient ionizing photons. To address these concerns, researchers have explored various factors such as the escape fraction (f\_esc) of ionizing photons from galaxies, the ionizing efficiency (xi\_ion), and the clumping factor (C) of the intergalactic medium [Davies2021, Park2022].
\section{Data and method}
In this study, we adopt a literature-anchored budget calculation approach to reconcile the reionization photon budget at z\~{}6. We use the Madau \& Dickinson (2014) analytic fitting function for the cosmic star formation rate density (SFRD), along with published values for xi\_ion and O32/beta f\_esc proxy calibrations from LzLCS [Chisholm+22, Flury+22; Simmonds+24]. Our method focuses on systematically reconciling these literature values to determine the required escape fraction of ionizing photons from star-forming galaxies.
\section{Result}
Our analysis reveals that reionization at z\~{}6 can be achieved if star-forming galaxies have an escape fraction f\_esc = 0.014 (+0.014/-0.007). This value is lower than the indirect-proxy-inferred f\_esc = 0.062 (+0.110/-0.039) derived from LzLCS O32/beta calibrations. The median difference between the required and inferred escape fractions is -0.045 dex-frac, with an 11\% probability of a shortfall in the photon budget. Importantly, this result holds under both O32 and beta calibrations.
\begin{figure}\centering\includegraphics[width=\columnwidth]{result.png}\caption{The reionization ionizing-photon-budget shortfall for star-forming galaxies is not robust to systematics at z\~{}6 (jwsthigh systematic corner)}\end{figure}
\section{Caveats}
However, it is essential to acknowledge the limitations of our approach. Our analysis relies on automated, single-selection, and uncalibrated measurements, which may introduce systematic uncertainties. The accuracy of our results depends heavily on the assumptions made in the literature-anchored budget calculation, including the choice of SFRD function and proxy calibrations. Furthermore, our study does not account for potential variations in xi\_ion and clumping factor C across different galaxy populations or environments. These factors may affect the actual photon budget and escape fraction requirements, emphasizing the need for further research to refine these estimates.
\section*{References}
\footnotesize
[Muoz2024] Muñoz et al. (2024). Reionization after JWST: a photon budget crisis?. bibcode 2024MNRAS.535L..37M \par [Davies2021] Davies et al. (2021). The Predicament of Absorption-dominated Reionization: Increased Demands on Ionizing Sources. bibcode 2021ApJ...918L..35D \par [Park2022] Park et al. (2022). Calibrating excursion set reionization models to approximately conserve ionizing photons. bibcode 2022MNRAS.517..192P \par [Duncan2015] Duncan et al. (2015). Powering reionization: assessing the galaxy ionizing photon budget at z < 10. bibcode 2015MNRAS.451.2030D \par [Madau2017] Madau (2017). Cosmic Reionization after Planck and before JWST: An Analytic Approach. bibcode 2017ApJ...851...50M

\end{document}
